\documentclass[11pt]{article}
\usepackage[margin=1in]{geometry}
\usepackage{amsmath,amssymb,amsthm,array,booktabs,tabularx}
\usepackage[T1]{fontenc}
\usepackage{lmodern}
\usepackage[colorlinks=true,urlcolor=blue,citecolor=blue,linkcolor=blue]{hyperref}
\newtheorem{theorem}{Theorem}
\newtheorem{lemma}[theorem]{Lemma}
\newtheorem{corollary}[theorem]{Corollary}
\newcommand{\Z}{\mathbb Z}
\newcommand{\N}{\mathbb N}
\newcommand{\card}[1]{\lvert #1\rvert}
\title{Shannon's circle-vertex conjecture and the\linebreak
Markey--Pfoertner seed-growth relation:\linebreak an expanding-hexagon proof}
\author{Haobo Ma\\\small ChronoAI Pte Ltd, Singapore
\and Wenlin Zhang\\\small National University of Singapore\\\small ChronoAI Pte Ltd, Singapore}
\date{16 November 2026}
\begin{document}
\maketitle
\begin{abstract}
We prove Shannon's three conjectured circle-vertex formulas in OEIS A374338,
FORMULA (revision 4, 5 July 2024; retained in revision 21), by identifying circle
intersections with two-neighbor growth on the triangular lattice
\cite{shannon}. The vertex count after $n\geq1$ steps is
$(3n^2+5n+c_r)/3$, where $r=n\bmod3$ and $(c_0,c_1,c_2)=(0,4,2)$.
A strip-clipped hexagon describes every collinear seed length, and becomes
an unclipped hexagon at the sharp threshold $n\geq\max(2,m-3)$ for $m$ seeds
and $n\geq2$. Its counts prove the three- and four-seed formulas recorded by Markey and
the conjectured five-seed difference relation observed by Pfoertner and
recorded by Markey. Corollaries give the periodic second differences and
the precise series identities linking Barker's circle-addition and petal
generating functions.
\end{abstract}

\section{The counts and their interpretation}
Let $a\ne b$ be points in the Euclidean plane, and put $R=\card{a-b}$.
Starting from $V_0=\{a,b\}$, retain all points and set
\[
 V_{n+1}=V_n\cup\{p:\ \card{p-u}=\card{p-v}=R
                  \text{ for distinct }u,v\in V_n\}.
\]
Thus $a(n)=\card{V_n}$ counts all retained points, not just the points newly
added: $a(0)=2$, $a(1)=4$, and $a(2)=8$.
This is the convention of the tabulated sequence in Shannon's entry
\cite[NAME, DATA, FORMULA]{shannon}.

For lattice growth, use axial coordinates $(x,y)\in\Z^2$ and the six directions
\[
 D=\{(1,0),(0,1),(-1,1),(-1,0),(0,-1),(1,-1)\}.
\]
For integers $m\geq2$, let $S_m(0)=\{(j,0):0\leq j<m\}$, and obtain
$S_m(n+1)$ by retaining $S_m(n)$ and adjoining each site with at least two
distinct neighbors in $S_m(n)$. Write $s_m(n)=\card{S_m(n)}$.
For $m=2,3,4,5$ these are the constructions in
\cite{shannon,three,four,five}. The use of \emph{at least two} is essential:
the exactly-two rule in A338097 describes a different process \cite{sigrist}.

\begin{theorem}\label{thm:counts}
For every pair of distinct starting points, $a(n)=s_2(n)$ for all $n\geq0$.
For $n\geq1$,
\begin{equation}\label{eq:shannon}
 a(n)=\frac{3n^2+5n+c_{n\bmod3}}{3},\qquad(c_0,c_1,c_2)=(0,4,2).
\end{equation}
For $n\geq0$ the three- and four-seed counts are
\begin{align}
 s_3(n)&=n^2+3n+
 \begin{cases}1,&n>0\text{ and }3\mid n,\\3,&\text{otherwise},\end{cases}\label{eq:three}\\
 s_4(n)&=\frac{3n^2+13n+12+2(n\bmod3)}{3}.\label{eq:four}
\end{align}
The five-seed count has $s_5(0)=5$, $s_5(1)=13$, and, for $n\geq2$,
\begin{equation}\label{eq:five}
 s_5(n)=\frac{3n^2+17n+25-2((n-1)\bmod3)}{3}
       =a(n)+4n+7.
\end{equation}
\end{theorem}

Shannon explicitly labelled \eqref{eq:shannon} conjectured in 2024.
Markey recorded \eqref{eq:three}--\eqref{eq:five} in October 2026; the entry
history identifies Pfoertner's observation behind the last equality in
\eqref{eq:five} \cite[4--5 October 2026 discussion]{fivehistory}.
Irvine recorded generating functions for the three-, four-, and five-seed sequences
in October 2026 \cite{three,four,five}. Earlier, Barker recorded a generating
function for the circle-addition sequence A264788 in December 2015
\cite{barker}. These algebraic formulas are credited separately from the
geometric identification and the all-step shape proof below. The common
mechanism is a congruence-filtered expanding hexagon. Retaining a horizontal
reachability strip extends that description from the four named sequences to
every seed length; counting it gives the formulas and the stabilization time.
The zero-index exception in the three-seed formula incorporates Xausa's
correction of 10 October 2026 \cite{threehistory}.
Neder's earlier periodic-diagram argument for A136289 starts with three
mutually touching pennies and counts new additions \cite{neder}; this is
a related predecessor, not the collinear-seed total enumerated here.

\section{Circle intersections stay in the triangular lattice}
Identify the plane with $\mathbb C$ by a similarity carrying $a$ to $0$ and
$b$ to $1$. Put $\omega=(1+i\sqrt3)/2$ and
$\Lambda=\{x+y\omega:(x,y)\in\Z^2\}$. Its squared norm is
$Q(x,y)=x^2+xy+y^2$.

\begin{lemma}\label{lem:circle}
If distinct $u,v\in\Lambda$ have intersecting unit circles, all their
intersection points belong to $\Lambda$ and are unit neighbors of both
centers. Consequently $V_n$ is the image of $S_2(n)$ under the initial
similarity for every $n$.
\end{lemma}
\begin{proof}
Translate $u$ to $0$ and write $v=x+y\omega$. Intersection implies
$0<Q(x,y)\leq4$. The identity
$4Q(x,y)=(2x+y)^2+3y^2$ restricts $y$ to $-2\leq y\leq2$.
For $y=0,1,2$ the possible $x$ are respectively
$\{-2,-1,1,2\}$, $\{-2,-1,0,1\}$, and $\{-2,-1,0\}$;
the negative-$y$ cases follow by negating both coordinates. These are
the six rotations of $1$, the six rotations of $1+\omega$, and the
six rotations of $2$, with squared norms $1$, $3$, and $4$.
Rotations by multiples of $\pi/3$ preserve $\Lambda$.
For centers $0,1$ the intersections are $\omega,\overline\omega$;
for centers $0,1+\omega$ they are $1,\omega$; for centers $0,2$
there is the single tangency point $1$. These lists follow by subtracting
the two circle equations and substituting in the first. Every listed point
is in $\Lambda$. The integral solutions of $Q=1$ are exactly the six
directions $D$. Thus a site is added by the circle rule precisely when it
has two distinct occupied lattice neighbors. Induction proves the last assertion.
\end{proof}

\section{A common expanding hexagon}
For integers $A,L$ define
\[
 H_{A,L}=\{(u,v)\in\Z^2:0\leq u,v\leq L,\quad
                         A\leq u+v\leq2L-A\}.
\]
The linear map $T(x,y)=(x-y,x+2y)$ sends $D$ to
\[
 M=\{(1,1),(-1,2),(2,-1),(-1,-1),(1,-2),(-2,1)\}
\]
and has image $\Gamma=\{(u,v)\in\Z^2:u\equiv v\pmod3\}$.
Its inverse on $\Gamma$ is $((2u+v)/3,(v-u)/3)$.

\begin{lemma}\label{lem:update}
Suppose $A\geq3$ and $L\geq A+2$. For $p=(u,v)\in\Z^2$,
\[
 p+(1,1)\in H_{A+1,L+2}
 \quad\Longleftrightarrow\quad
 p\in H_{A,L}\ \text{or}\ p+d,p+e\in H_{A,L}
 \text{ for distinct }d,e\in M.
\]
\end{lemma}
\begin{proof}
The left side is equivalent to
\begin{equation}\label{eq:expanded}
 -1\leq u,v\leq L+1,\qquad A-1\leq u+v\leq2L-A+1.
\end{equation}
Old sites satisfy these bounds. If two distinct neighbors belong to
$H_{A,L}$, each of the bounds follows from the list $M$: only one move
has first coordinate $2$, only one has first coordinate $-2$, and likewise
for the second coordinate; only one has coordinate sum $2$, and only one
has coordinate sum $-2$. Thus two distinct moves cannot violate any bound
in \eqref{eq:expanded}.

Conversely, assume \eqref{eq:expanded} and $p\notin H_{A,L}$.
Use the first applicable row of the following table. In the last two
columns the first pair is used under the stated condition, and the second
pair otherwise. Each pair gives $d,e\in M$ with $p+d,p+e\in H_{A,L}$.
\begin{center}\small
\begin{tabular}{@{}llll@{}}
\toprule
Boundary & Condition & First pair & Second pair\\
\midrule
$u+v=A-1$ & $u\geq1$ & $(1,1),(-1,2)$ & $(1,1),(2,-1)$\\
$u+v=2L-A+1$ & $u\leq L-1$ & $(-1,-1),(1,-2)$ & $(-1,-1),(-2,1)$\\
$u=-1$ & $v\leq A+1$ & $(2,-1),(1,1)$ & $(2,-1),(1,-2)$\\
$u=L+1$ & $L-v\leq A+1$ & $(-2,1),(-1,-1)$ & $(-2,1),(-1,2)$\\
$v=-1$ & $u\leq A+1$ & $(-1,2),(1,1)$ & $(-1,2),(-2,1)$\\
$v=L+1$ & $L-u\leq A+1$ & $(1,-2),(-1,-1)$ & $(1,-2),(2,-1)$\\
\bottomrule
\end{tabular}
\end{center}
Here is an explicit check of the first row. If $u\geq1$, then
$v=A-1-u\geq-1$ gives $u\leq A\leq L-2$. The shifted first coordinates
$u+1,u-1$ and second coordinates $v+1,v+2$ lie in $[0,L]$;
their sums are $A+1,A$, both at most $2L-A$ since $L\geq A+2$.
If $u\leq0$, the same bounds hold for $(1,1),(2,-1)$, using
$-1\leq u\leq0$ and $A-1\leq v\leq A$.
The second row is the reflection $(u,v)\mapsto(L-u,L-v)$ of this check.
For the third row the two diagonal boundaries have already been excluded,
so $A+1\leq v\leq L+1$. If $v\leq A+1$, take $v=A+1$ and the indicated
pair. Otherwise $A+2\leq v\leq L+1$, and the second pair has first
coordinates $1,0$, second coordinates $v-1,v-2$, and sums $v,v-2$,
all within the old bounds. The fourth row follows by reflection, and
the last two by swapping coordinates. If none of the six boundaries occurs,
the integer inequalities put $p$ in $H_{A,L}$, a contradiction.
\end{proof}

\begin{theorem}\label{thm:shape}
For every $m\geq2$ and $n\geq2$, put $L=m-1+2n$. Then
\begin{equation}\label{eq:shape}
 T(S_m(n))+(n,n)=
 H_{n+1,L}\cap\Gamma\cap\{(u,v):|v-u|\leq3n\}.
\end{equation}
For $n\geq2$ the strip can be omitted exactly when $n\geq m-3$.
\end{theorem}
\begin{proof}
At step one the occupied rows are
$y=-1$, $1\leq x\leq m-1$;
$y=0$, $0\leq x\leq m-1$; and
$y=1$, $0\leq x\leq m-2$.
Taking two-neighbor growth once more gives the following complete rows:
\[
\begin{array}{c|ccccc}
y&-2&-1&0&1&2\\\hline
x&[2,m-1]&[0,m]&[0,m-1]&[-1,m-1]&[0,m-3].
\end{array}
\]
Intervals here contain integers, and an interval with lower endpoint larger
than upper endpoint is empty. On row $y=-2$, the only possible neighbors
are $(x,-1),(x-1,-1)$, so both are occupied exactly when
$2\leq x\leq m-1$. On row $y=2$ the corresponding pair
$(x,1),(x+1,1)$ gives $0\leq x\leq m-3$.
Row $y=-1$ acquires exactly the endpoints $x=0,m$ of the displayed interval:
each has an occupied horizontal neighbor and an occupied neighbor on row zero;
farther sites have at most one occupied neighbor. The reflection
$(x,y)\mapsto(x+y,-y)$ gives row $y=1$. Row zero retains its old interval;
the two next exterior sites each have just one occupied neighbor, and all
farther sites have none.
No row with $\card y>2$ can have an occupied neighbor at step one.
On the right of \eqref{eq:shape} at $n=2$, the strip gives $\card y\leq2$.
Substituting $u=x-y+2$, $v=x+2y+2$ gives exactly the five displayed
intervals. This proves the base case for every $m\geq2$.

For the induction put $A=n+1$ and let $p=(u,v)$ be the coordinates before
the translation by $(1,1)$ at the next step. The old hexagon is $H_{A,L}$.
Since $n+1\geq3$ and $m-1+2n\geq(n+1)+2$, Lemma~\ref{lem:update}
applies. Its bounds give the next hexagon, and every move changes $v-u$
by $0$ or $\pm3$, giving the next strip. Conversely, consider a point in
the next hexagon and strip. Both the moves and translation preserve $\Gamma$.
If $|v-u|\leq3n$, a point in the old hexagon is retained; for a point outside
it use the witness table.
On its diagonal boundaries the expanded coordinate bounds give
$|v-u|\leq n+2<3n$ (reflect for the upper boundary). Congruence then gives
$|v-u|\leq3n-3$, so both witnesses stay in the old strip. On $u=-1$ the
witnesses decrease or preserve the positive difference $v-u$; on $u=L+1$
they increase or preserve its negative value. The $v$ boundaries follow
by swapping coordinates. These changes cannot cross the opposite strip
boundary, since $n\geq2$.

The remaining possibility, by congruence, is $|v-u|=3n+3$.
For $v-u=3n+3$ use the two moves $(1,-2),(2,-1)$. Their differences are
$3n$, and the bounds $-1\leq u,v\leq L+1$ put their coordinates in $[0,L]$.
Moreover $3n+1\leq u+v\leq2L-3n-1$, so their sums $u+v-1,u+v+1$
belong to $[n+1,2L-n-1]$. Thus both are old occupied sites.
The negative case follows by swapping coordinates. This proves the induction.

If $n\geq m-3$, then $L\leq3n+2$; since $v-u$ is divisible by $3$,
the square already implies the strip. If instead $2\leq n<m-3$,
the point $(0,3n+3)$ lies in $H_{n+1,L}\cap\Gamma$ but not the strip.
Thus the threshold is exact.
\end{proof}

\section{Counting and consequences}
Let $C(n)$ count nonnegative integer pairs $(u,v)\in\N^2$ with $u+v\leq n$ and
$u\equiv v\pmod3$. In the square $0\leq u,v\leq L$ the congruence
condition gives
\[
 Q_L=\sum_{r=0}^2\left\lfloor\frac{L+3-r}{3}\right\rfloor^2.
\]
When the strip is redundant, the two excluded corners in \eqref{eq:shape}
each contain $C(n)$ points:
the upper corner is carried to the lower one by $(u,v)\mapsto(L-u,L-v)$.
They are disjoint because $n<L$. Hence the strongest common count is
\begin{equation}\label{eq:common}
 s_m(n)=Q_{m-1+2n}-2C(n)\qquad(m\geq2,\ n\geq\max(2,m-3)).
\end{equation}
For completeness the antidiagonal counts, and their sums, are
\[
\begin{array}{c|ccc}
n&3k&3k+1&3k+2\\\hline
\card{\{u+v=n,\ u\equiv v\pmod3\}}&k+1&k&k+1\\
2C(n)&3k^2+3k+2&3k^2+5k+2&3k^2+7k+4.
\end{array}
\]
Indeed the first row enumerates $u$ in the required residue class modulo
$3$. Summing successive triples gives the second row, starting with
$C(0)=1$. Substitution in \eqref{eq:common} gives
\eqref{eq:shannon}--\eqref{eq:five}. For example, at $m=2$ the three
residue classes give $9k^2+5k$, $9k^2+11k+4$, and $9k^2+17k+8$.
Increasing $L$ by $3$ increases every square residue count by $1$;
thus $Q_{L+3}-Q_L=2(L+1)+3=2L+5$. Taking $L=2n+1$ gives
$s_5(n)-s_2(n)=4n+7$. The directly counted values
$s_m(0)=m$, $s_m(1)=3m-2$ give all remaining boundary cases.
Together with Lemma~\ref{lem:circle}, this proves Theorem~\ref{thm:counts}.

\begin{corollary}\label{cor:seedlength}
For $m\geq2$ and $n\geq\max(2,m)$,
\[
 s_{m+3}(n)-s_m(n)=4n+2m+3.
\]
For fixed $m$, the eventual count is a quadratic quasi-polynomial of period
dividing three, with leading coefficient one. The unclipped hexagon formula
for every seed length already at step two is not asserted in the posing
sources; it fails for six seeds: the hexagon contains $30$ sites but growth
has $28$, omitting axial sites $(1,3)$ and $(4,-3)$.
\end{corollary}
\begin{proof}
Both counts have stabilized to \eqref{eq:common}; subtract and use
$Q_{L+3}-Q_L=2L+5$ with $L=m-1+2n$.
The residue formulas for $Q_L$ and $2C(n)$ are quadratic, and their leading
terms give $s_m(n)=n^2+O_m(n)$. The two extra sites in the six-seed example
have $|y|=3$, so are excluded by \eqref{eq:shape}; the step-two row table
counts the remaining $28$ sites.
\end{proof}

\begin{corollary}\label{cor:second}
Markey's second-difference pattern in the A400701 history (revision 4,
4 October 2026, COMMENTS) holds for every $n\geq4$:
\[
 s_5(n)-2s_5(n-1)+s_5(n-2)=
 \begin{cases}2,&n\equiv0\pmod3,\\4,&n\equiv1\pmod3,\\0,&n\equiv2\pmod3.
 \end{cases}
\]
\end{corollary}
\begin{proof}
For the three indices the linear difference $s_5(j)-a(j)=4j+7$
cancels. Subtracting the successive increments
$a(j)-a(j-1)=2j+2\mathbf1_{j\equiv1\ (3)}$ gives the displayed cases.
\end{proof}

\begin{corollary}\label{cor:gf}
For $2\leq m\leq5$, put $F_m(z)=\sum_{n\geq0}s_m(n)z^n$ and
$P(z)=(1-z)^2(1-z^3)$. Then $F_m(z)=N_m(z)/P(z)$, where
\[
\begin{array}{c|l}
m&N_m(z)\\\hline
2&2+2z^2+4z^4-2z^5\\
3&3+z+2z^2-3z^3+5z^4-2z^5\\
4&4+2z+2z^2-4z^3+2z^4\\
5&5+3z+2z^2-5z^3+z^4-2z^5+2z^6.
\end{array}
\]
In particular $s_m(n)=2s_m(n-1)-s_m(n-2)+s_m(n-3)
-2s_m(n-4)+s_m(n-5)$ for $n\geq7$ and $2\leq m\leq5$.
\end{corollary}
\begin{proof}
For $m=2$, \eqref{eq:shannon} gives
$s_2(n)-s_2(n-1)=2n+2\,\mathbf1_{n\equiv1\ (3)}$ for $n\geq2$,
with the exceptional first increment $2$.
For $m=3,4,5$ the formulas give a linear function plus a periodic
function for the first differences when both indices are at least $2$.
Subtracting such differences three steps apart is constant, and taking
one more difference gives zero. Equivalently $P(z)F_m(z)$ has no
coefficients past degree $6$. Multiplying its first seven coefficients,
using Theorem~\ref{thm:counts}, gives the displayed numerators.
Since $P(z)=1-2z+z^2-z^3+2z^4-z^5$, coefficient comparison gives the recurrence.
\end{proof}

The numerators for $m=3,4,5$ agree with the generating functions already
recorded in the corresponding OEIS entries \cite{three,four,five}; the
shape argument supplies their geometric justification. If $b(j)$ is the
sequence specified by Barker's generating function in A264788, then
\begin{equation}\label{eq:barker}
 \sum_{j\geq0}b(j)z^j=2+2z^2+z(1-z)F_2(z)
 =\frac{2(1+z^2-z^3+z^4+z^5)}{(1-z)(1-z^3)}.
\end{equation}
This follows by substituting $N_2/P$ and simplifying, and equals Barker's
recorded expression because $(1-z)(1-z^3)=(1-z)^2(1+z+z^2)$.
Therefore $a(n)-a(n-1)=b(n+1)$ for $n\geq2$.
The phrase ``first differences'' in the A374338 cross-reference requires
this indexing qualification: $a(3)-a(2)=6$, whereas $b(3)=4$.
Equation~\eqref{eq:barker} is a numerical generating-function identity;
the distinct wording of the geometric construction of A264788 is not
identified with the growth rule here.

\begin{corollary}\label{cor:petal}
Let $c(j)$ be the coefficients of Barker's recorded petal series
\cite[FORMULA, 13 December 2015]{petals}, namely
\[
 C(z)=\sum_{j\geq0}c(j)z^j
 =\frac{z(1+5z+6z^2+11z^3-5z^4)}{(1-z)^2(1-z^3)}.
\]
Then $c(n+3)=3s_5(n)-2$ for every $n\geq2$.
Also $c(j)-2c(j-1)+c(j-2)$ equals $6,12,0$ for
$j\equiv0,1,2\pmod3$, respectively, for all $j\geq3$.
\end{corollary}
\begin{proof}
Substituting the numerator $N_5$ of Corollary~\ref{cor:gf} gives
\[
 C(z)=z+7z^2+6z^3+6z^4
       +z^3\left(3F_5(z)-\frac{2}{1-z}\right).
\]
The polynomial has degree four, proving the first assertion.
For the second, multiply $C(z)$ by $(1-z)^2$ and compare the coefficients
of $(z+5z^2+6z^3+11z^4-5z^5)/(1-z^3)$ from index three onward.
\end{proof}
This proves the series form of Pfoertner's empirical observation in the
A400701 discussion of 4 October 2026, 18:38 EDT \cite{fivehistory}, and the
series form of the cyclic second differences noted by Ngaokrajang
\cite[COMMENTS]{petals}. Here $c$ is defined by Barker's displayed series;
identifying that series independently with geometric petal counts would
require a separate petal-counting argument.

\section{Coverage and verification}
\begin{center}\small
\begin{tabularx}{\textwidth}{@{}lX@{}}
\toprule
Source item & Coverage\\\midrule
A374338, three conjectured formulas & Proved for every $n\geq1$ and every distinct starting pair.\\
A400699, A400700, FORMULA & Recorded three- and four-seed counts proved for all $n\geq0$.\\
A400701, FORMULA & Five-seed formula and conjectured difference relation proved for $n\geq2$; both initial terms counted separately.\\
A400701, revision 4 COMMENTS & Proved: Markey's observed periodic second differences, for all $n\geq4$ (Corollary~\ref{cor:second}).\\
A400701, 4 October discussion & Pfoertner's petal relation proved for the coefficients of Barker's cited series; independent geometric petal enumeration is not covered (Corollary~\ref{cor:petal}).\\
A264788, FORMULA and cross-reference & Barker's expression cited; shifted difference identity proved in \eqref{eq:barker}.\\
A374337, A374339 \cite{regions} & Region and curved-edge counts are not determined here. Their Euler relation alone does not separate the two counts.\\
A338097, A008458 \cite{sigrist,shells} & Respectively a different exactly-two update and a distance-shell coordination sequence; neither count is asserted equal to $s_2(n)$.\\
A265676 \cite{petals} & Prior formula cited (Barker, 2015); its coefficient second differences and five-seed relation proved. Geometric petal counts are not covered.\\
A136289 \cite{neder} & Prior proof cited (Neder, 2018) for its periodic diagram; its noncollinear-seed, new-addition problem is not reproved here.\\
Arbitrary collinear seed lengths & Not posed; proved: strip-clipped shape, sharp eventual hexagon threshold, and width-three count relation (Theorem~\ref{thm:shape}, Corollary~\ref{cor:seedlength}).\\
A359569, A371373, A371254 \cite{othercircles} & Variable radii or regular-polygon arc constructions have different initial data or geometry; their questions are not covered.\\
A400975, A400002, A400977 \cite{otherrules} & Exactly-two or adjacent-pair update variants are not covered.\\
\bottomrule
\end{tabularx}
\end{center}

The written proof covers every collinear seed length. The following pinned Lean
declarations specify the formal scope for the two- and five-seed cases:
\begin{itemize}
\item \href{https://github.com/the-omega-institute/trureturing/blob/2f8421f6b90d37b4815b887f6c60987553831b20/D5/S3/Geometry/ShannonFixedRadiusCircleVertices.lean}{\texttt{ShannonFixedRadiusCircleVertices.lean}}:
\texttt{ShannonA374338.result} proves all three formulas for arbitrary
distinct Euclidean starting points. The finite-shape count difference is
proved by \texttt{count\_difference} in \texttt{HexShapeCounts}.
\item \href{https://github.com/the-omega-institute/trureturing/blob/2f8421f6b90d37b4815b887f6c60987553831b20/D5/S3/StatisticalMechanics/Percolation/TriangularTwoNeighborGrowth.lean}{\texttt{TriangularTwoNeighborGrowth.lean}}:
\texttt{AxialGrowth.exact\_shape} gives the shape for $m=2,5$ and $n\geq2$;
\texttt{CircleGrowth.intersection\_lattice} gives the Euclidean-to-lattice step.
\item \href{https://github.com/the-omega-institute/trureturing/blob/2f8421f6b90d37b4815b887f6c60987553831b20/D5/S3/StatisticalMechanics/Percolation/MarkeyFiveSiteGrowthDifference.lean}{\texttt{MarkeyFiveSiteGrowthDifference.lean}}:
\texttt{MarkeyA400701.result} proves the five-seed lattice count equals
the two-point circle count plus $4n+7$ for $n>1$.
\end{itemize}
The arbitrary-seed extension, three- and four-seed statements, periodic
differences and generating-function consequences are written proofs above;
the linked declarations have the narrower scopes just specified.
The accompanying \texttt{verify.py} checks finite
growth states, boundary witnesses, and exact-integer coefficient identities;
its ranges and results are given in \texttt{verification.json}.

\begin{thebibliography}{14}
\bibitem{shannon} S. R. Shannon, OEIS A374338, \emph{Fixed-radius circle vertices},
5 July 2024; revision 21, 10 November 2026, NAME, DATA, FORMULA and CROSSREFS.
\url{https://oeis.org/A374338}.
Original conjectured formulas: revision 4 in
\url{https://oeis.org/history?seq=A374338\&start=10}.
\bibitem{three} A. Markey, OEIS A400699, \emph{Three adjacent collinear seeds},
3 October 2026; revision 38, 3 November 2026, FORMULA.
Generating function credited there to S. A. Irvine, 9 October 2026.
\url{https://oeis.org/A400699}.
\bibitem{four} A. Markey, OEIS A400700, \emph{Four adjacent collinear seeds},
3 October 2026; revision 35, 3 November 2026, FORMULA.
Generating function credited there to S. A. Irvine, 9 October 2026.
\url{https://oeis.org/A400700}.
\bibitem{five} A. Markey, OEIS A400701, \emph{Five adjacent collinear seeds},
3 October 2026; revision 31, 25 October 2026, FORMULA.
Generating function credited there to S. A. Irvine, 19 October 2026.
\url{https://oeis.org/A400701}.
\bibitem{fivehistory} A. Markey and H. Pfoertner, OEIS A400701 edit history,
4--5 October 2026, revisions 4 and 8--11 and accompanying discussion.
\url{https://oeis.org/history?seq=A400701\&start=20}.
\bibitem{barker} K. Ngaokrajang, OEIS A264788, \emph{Circle additions},
25 November 2015; revision 20, 16 February 2025, FORMULA.
Generating function credited there to C. Barker, 10 December 2015.
\url{https://oeis.org/A264788}.
\bibitem{regions} S. R. Shannon, OEIS A374337 and A374339,
\emph{Regions and curved edges in the circle construction}, 5 July 2024.
\url{https://oeis.org/A374337}; \url{https://oeis.org/A374339}.
\bibitem{sigrist} R. Sigrist, OEIS A338097, \emph{Exactly-two-neighbor growth},
10 October 2020, NAME. \url{https://oeis.org/A338097}.
\bibitem{shells} N. J. A. Sloane, OEIS A008458, \emph{Coordination sequence for
the triangular lattice}, NAME. \url{https://oeis.org/A008458}.
\bibitem{petals} K. Ngaokrajang, OEIS A265676, \emph{Flower-of-Life petals},
13 December 2015, COMMENT and FORMULA. Recurrence and generating function
credited there to C. Barker, 13 December 2015. \url{https://oeis.org/A265676}.
\bibitem{othercircles} S. R. Shannon, OEIS A359569 (6 January 2023),
A371373 (20 March 2024), and A371254 (16 March 2024), NAME and COMMENTS.
\url{https://oeis.org/A359569}; \url{https://oeis.org/A371373};
\url{https://oeis.org/A371254}.
\bibitem{otherrules} A. Markey, OEIS A400975 and A400002 (11 October 2026),
and A400977 (18 October 2026), NAME and COMMENTS.
\url{https://oeis.org/A400975}; \url{https://oeis.org/A400002};
\url{https://oeis.org/A400977}.
\bibitem{threehistory} P. Xausa, OEIS A400699 edit history, revision 28,
10 October 2026, 14:25 EDT, FORMULA and discussion: correction of the
zero-index case. \url{https://oeis.org/history?seq=A400699\&start=10}.
\bibitem{neder} C. Mallows, OEIS A136289, 13 April 2008; C. Neder,
COMMENTS, 22 June 2018, periodic-triangle argument for the recorded formulas.
\url{https://oeis.org/A136289}.
\end{thebibliography}
\noindent\small Copyright \copyright\ 2026 Haobo Ma and Wenlin Zhang.
Licensed under Creative Commons Attribution 4.0 International (CC BY 4.0).
\url{https://creativecommons.org/licenses/by/4.0/}.
\end{document}
